\documentclass[
	% -- opções da classe memoir --
	12pt,				% tamanho da fonte
	openright,			% capítulos começam em pág ímpar (insere página vazia caso preciso)
	oneside,			% para impressão em verso e anverso. Oposto a oneside
	a4paper,			% tamanho do papel. 
	% -- opções da classe abntex2 --
	chapter=TITLE,		% títulos de capítulos convertidos em letras maiúsculas
	%section=TITLE,		% títulos de seções convertidos em letras maiúsculas
	subsection=TITLE,	% títulos de subseções convertidos em letras maiúsculas
	%subsubsection=TITLE,% títulos de subsubseções convertidos em letras maiúsculas
	% -- opções do pacote babel --
	english,			% idioma adicional para hifenização
	brazil,				% o último idioma é o principal do documento
	]{abntex2}

% ---
% PACOTES
% ---

% ---
% Pacotes fundamentais 
% ---
\usepackage{lmodern}			% Usa a fonte Latin Modern
\usepackage[T1]{fontenc}		% Selecao de codigos de fonte.
\usepackage[utf8]{inputenc}		% Codificacao do documento (conversão automática dos acentos)
\usepackage{indentfirst}		% Indenta o primeiro parágrafo de cada seção.
\usepackage{color}				% Controle das cores
\usepackage{graphicx}			% Inclusão de gráficos
\usepackage{microtype} 			% para melhorias de justificação

% ---

%--------------Pacotes de citações
\usepackage[brazilian,hyperpageref]{backref}	 % Paginas com as citações na bibl
\usepackage[alf]{abntex2cite}	% Citações padrão ABNT

\usepackage{hyperref}

% --- 
% CONFIGURAÇÕES DE PACOTES
% --- 

% ---
% Configurações do pacote backref
% Usado sem a opção hyperpageref de backref
\renewcommand{\backrefpagesname}{Citado na(s) página(s):~}
% Texto padrão antes do número das páginas
\renewcommand{\backref}{}
% Define os textos da citação
\renewcommand*{\backrefalt}[4]{
	\ifcase #1 %
		Nenhuma citação no texto.%
	\or
		Citado na página #2.%
	\else
		Citado #1 vezes nas páginas #2.%
	\fi}%
% ---

% alterando o aspecto da cor azul
\definecolor{blue}{RGB}{41,5,195}

% informações do PDF
\makeatletter
\hypersetup{
     	%pagebackref=true,
		pdftitle={\@title}, 
		pdfauthor={\@author},
    	pdfsubject={\imprimirpreambulo},
	    pdfcreator={LaTeX with abnTeX2},
		pdfkeywords={abnt}{latex}{abntex}{abntex2}{plano de trabalho}, 
		colorlinks=true,       		% false: boxed links; true: colored links
    	linkcolor=blue,          	% color of internal links
    	citecolor=blue,        		% color of links to bibliography
    	filecolor=magenta,      		% color of file links
		urlcolor=blue,
		bookmarksdepth=4
}
\makeatother
% --- 

% --- 
% Espaçamentos entre linhas e parágrafos 
% --- 

% O tamanho do parágrafo é dado por:
\setlength{\parindent}{1.3cm}

% Controle do espaçamento entre um parágrafo e outro:
\setlength{\parskip}{0.2cm}  % tente também \onelineskip

\usepackage{fancyhdr}
\usepackage{lipsum} % Apenas para gerar texto fictício

% Configuração do cabeçalho especial para a primeira página
\fancypagestyle{primeirapagina}{
    \fancyhf{} % Limpa todos os campos de cabeçalho e rodapé
    \renewcommand{\headrulewidth}{0pt} % Remove a linha horizontal do topo
    
    % Define as logos alinhadas na vertical por sua base ([b]) ou centro ([c])
    \fancyhead[L]{\includegraphics[height=1.5cm]{Images/logo-ufu.png}}
    \fancyhead[R]{\includegraphics[height=1.2cm]{Images/logo-facom.png}}
    
    % Ajusta a distância do cabeçalho para o texto nesta página
    \setlength{\headsep}{1,5cm} 
}

% ---
% compila o indice
% ---
\makeindex
% ---

%-----------------------------------------
\begin{document}

% Ativa o cabeçalho das logos apenas nesta primeira página
\thispagestyle{primeirapagina}

% Retira espaço extra obsoleto entre as frases.
\frenchspacing 

\textbf{Disciplina:} Selecionar a opção correta conforme sua matrícula
%\textbf{Disciplina:} FACOM31701 - Trabalho de Conclusão de Curso 1 %BSI-UDI
%\textbf{Disciplina:} FACOM32703 - Projeto de Graduação I - %BSI-UDI Nova
%\textbf{Disciplina:} FACOM31703 - Trabalho de Conclusão de Curso 1 %BSI-MC (PPC-2016)
%\textbf{Disciplina:} FACOM33701 - Trabalho de Conclusão de Curso 1 %BSI-MC (PPC-2022)
%\textbf{Disciplina:} GBC072 - Projeto de Graduação 1 - %BCC
%\textbf{Disciplina:} FACOM31704 - Metodologia Científica em Computação e Informática %%BCC Nova

\textbf{Professor:} Selecionar a alternativa conforme o professor responsável pela disciplina 
%\textbf{Professor:} Prof. Dr. Rafael Dias Araújo
%\textbf{Professor:} Prof. Dr. Victor Sobreira

\textbf{Estudante:} Nome Completo

\textbf{Orientador(a):} Prof.(a) Nome Completo do Professor(a) (se, houver, indicar co-orientador abaixo)

%\textbf{Coorientador(a):} Prof.(a) Nome Completo do(a) Professor(a) Coorientador

\textbf{Tipo de trabalho: } Selecionar conforme modelo sugerido
%\textbf{Tipo de trabalho: } Projeto de Pesquisa Tradicional OU Relatório Técnico Profissional OU Desenvolvimento de Software OU Plano de Negócios


% ------------------ Título -----------------------
\section*{\centerline{\Large \textbf{Entrega 3 (E3) - Trabalhos relacionados}}}
% -------------------------------------------------

%------------------- Corpo de dados --------------
\vspace*{0.5cm}

\textbf{1) Identifique, juntamente com o seu orientador(a) e com o auxílio de mecanismos de pesquisa e bibliotecas digitais (ex. Google Acadêmico, IEEEXplore, ACM Digital Library, ...), quais os principais trabalhos relacionados do seu projeto de TCC/PG (entre 3 e 5 trabalhos). Garanta que a bibliografia elaborada no arquivo ref.bib contenha todos os metadados necessários para identificação do veículo de publicação dos trabalhos relacionados.}

\textbf{2) Faça uma leitura dinâmica e escreva um parágrafo sobre cada um dos trabalhos relacionados identificados no item "1". Fale sobre o objetivo do trabalho, o método, os principais resultados, como ele se relaciona com o seu projeto.  Identifique o nível de impacto de cada um dos trabalhos com base no seu veículo de publicação (conferência ou revista científica) e outros indicadores de relevância (ex. Qualis, citações,  índice Scopus, produtos derivados, patentes, ...).}

\textbf{3) Ao final do texto, apresente uma tabela comparativa sobre as principais características dos trabalhos relacionados.}

\textbf{4) Em seguida, estude e elabore um resumo sobre o Capítulo 7 sobre Revisão Bibliográfica do livro de referência \cite{metodologiaWazlawick}, deixando claro os principais pontos de aplicação do estudo em seu trabalho.}

\textbf{5) Feche o texto declarando o uso ou não de Inteligência Artificial para produção deste trabalho.}

Observações:
\begin{enumerate}
    \item Leia outras monografias para obter uma compreensão mais clara sobre como os trabalhos relacionados devem ser descritos. 
    \item Os parágrafos sobre os trabalhos relacionados devem apresentar um mínimo de coesão. Existem diversos conectivos que podem ser utilizados para auxiliar nesta tarefa.
    \item No caso do \textbf{relatório técnico profissional}, faça um levantamento dos trabalhos relacionados ao seu tema (com base nas orientações iniciais), complementando com questões técnicas tais como, tecnologias utilizadas na execução do trabalho, estratégias de integração dos sistemas, modelo de segurança de tecnologia da informação adotado e processos de engenharia de software.
    \item No caso do \textbf{desenvolvimento de um software}, faça um levantamento dos trabalhos relacionados ao seu tema  (com base nas orientações iniciais), complementando com aspectos técnicos sobre o desenvolvimento dos sistemas, incluindo os requisitos funcionais do sistema, diagrama de casos de uso e detalhamento dos requisitos não funcionais  do sistema. 
    \item No caso de um \textbf{plano de negócios}, faça um levantamento dos trabalhos relacionados ao seu tema (com base nas orientações iniciais), complementando com o indicativo sobre os principais concorrentes do produto e/ou serviço especificado no plano de negócios, descritivo das propostas, levantamento de  mercado concorrente, incluindo preços aplicados e características diferenciais. 
    \item É de suma importância que todos os parágrafos escritos sejam coesos. Ou seja, o texto como um todo precisa fazer sentido;
    \item Use o formato de citações do Latex para apresentar as referências em todo o texto. Deve ser entendido como as citações funcionam no Latex, em particular com o BibTeX. 
    \item O posicionamento de citações no texto e orientações da ABNT serão discutidos nas aulas. Lembre-se que no modelo fornecido para a entrega 1, além do arquivo \texttt{main.tex} utilizado para redigir seu texto, há um arquivo chamado \texttt{ref.bib} com alguns exemplos de referências. Um tutorial sobre o overleaf está disponível em \cite{TutorialOverleaf}. 
    \item Junto da entrega, o aluno deverá anexar um comprovante de concordância do seu orientador quanto ao texto sendo submetido. O comprovante de concordância poderá ser um e-mail do seu orientador, onde ele manifesta autorização para submeter o texto, ou uma assinatura eletrônica no PDF submetido feita pelo orientador, preferencialmente pelo sistema \href{https://www.gov.br/governodigital/pt-br/identidade/assinatura-eletronica}{GovBr} (ou outra plataforma de assinaturas conhecida e com validação pública).
\end{enumerate}

\chapter{Trabalhos Relacionados}

Parágrafo sobre o Trabalho 1...

Parágrafo sobre o Trabalho 2...

Parágrafo sobre o Trabalho 3...

Parágrafo sobre o Trabalho 4...

Parágrafo sobre o Trabalho 5...

Exemplo de Tabela em LaTeX:
\begin{center}
\begin{tabular}{ |c|c|c| } 
 \hline
 cell1 & cell2 & cell3 \\ 
 cell4 & cell5 & cell6 \\ 
 cell7 & cell8 & cell9 \\ 
 \hline
\end{tabular}
\end{center}


\section{Resumo sobre o Capítulo 7}

Resumo sobre o Capítulo 7 de Revisão Bibliográfica do livro de referência \cite{metodologiaWazlawick}.

\section{Declaração sobre o uso de Inteligência Artificial}

Inclua um parágrafo declarando o uso ou não de Inteligência Artificial para produção deste trabalho. A declaração deve conter a(s) ferramenta(s) utilizadas e como/onde foram utilizadas.

\bibliography{ref}

\end{document}
